\documentclass[12pt]{article}

\addtolength{\textwidth}{.8in}
\addtolength{\hoffset}{-.4in}
\addtolength{\textheight}{1in}
\addtolength{\voffset}{-1in}

\usepackage{amsmath}
\usepackage{hyperref}
\usepackage{graphicx}
\usepackage{upquote} % to convert funny quotes to straight quotes
\usepackage{amsfonts}
\usepackage{paralist}
\setlength{\parindent}{15pt}

\renewcommand*{\thefootnote}{\fnsymbol{footnote}}


\begin{document}

\begin{center}
\begin{Large}
DATA 219 Notes -- Apr. 3, 2018
\end{Large}
\end{center}

\section{Connecting to an external API}

I called the screen-scraping scenario ``the worst of all possible worlds."
Using an API, the subject of today's lesson, is sort of the best. I say ``sort
of" because it's certainly not the \textit{easiest}. The easiest is when you
have a ready-made CSV file with exactly the data you want, and you suck it
into a DataFrame, done. The thing is, much of the data you'll be interested in
analyzing as a Data Scientist wasn't put together like that for you. Instead,
it exists as part of a vast database somewhere, and you need flexibility and
power to extract on demand exactly the parts you want. This is what a good API
lets you do.

API actually stands for ``Application Programming Interface," which is a
pretty uninformative acronym. What it really means is ``a set of functions
someone has provided for you to call in order to access their data, together
with clean documentation telling you precisely what each of those functions
do." As you'll see with experience, some APIs are better designed than others,
and make it easier to get what you want. In this module we'll look at two
examples and generalize some common principles from them.

\subsection{REST APIs}

Another strange pseudo-acronym you'll encounter is ``REST"\footnote{This sort
of stands for ``REpresentational State Transfer," another utterly unhelpful
phrase.}, which is a \textit{style} of API. REST APIs have the following
characteristics:

\begin{itemize}
\itemsep.1em

\item They're designed to be used over the Web. (Although ``over the Web" may
not mean what you think it means; for one, it \textit{doesn't} mean ``using a
browser." Instead, it means ``accessing the data using HTTP requests in the
form of structured URLs/hyperlinks." We'll be writing \textit{Python code} to
retrieve what appear to be pseudo-Web-pages outside our browser.)

\item They're simple and elegant. Unlike their predecessor (and competitor)
protocols SOAP and RPC, REST was designed and promoted by
people\footnote{Originally Roy Fielding in his 2000 Ph.D. dissertation from
UC-Irvine.} who didn't like clutter and who wanted communication to be
streamlined and compactly specified. REST APIs are typically nice to use.

\item They're based directly on the HTTP protocol itself, and under the hood
make use of the HTTP operations like \texttt{GET}, \texttt{HEAD},
\texttt{POST}, \texttt{DELETE}, \texttt{PUT}, \textit{etc.} Because of this,
different REST APIs tend to be pretty uniform in how they work.

\item They're ``stateless," as HTTP is itself. This means the server retains
no \textit{memory} of the past requests from a client. (``Server" is a word
meaning ``a machine on the Internet that you want to access," and ``client"
here essentially means ``your own computer, which makes requests of the
server.") This has some limitations, but makes APIs quick and simple to use.

\item The data they provide can come in a variety of formats, typically JSON
or XML.

\end{itemize}

\subsubsection{Other terms}

APIs that conform to Roy Fielding's original spec are known colloquially as
``RESTful." 

You'll also sometimes hear the buzzwords ``Web services" or ``SOA"
(Service-oriented Architecture). Although these terms are a bit broader and
communicate something subtly different than just ``the use of Web-based APIs"
(REST or otherwise), in practice you will see terms like this used jointly to
refer to the same technologies. If you hear someone talking about ``a Web
services approach," your ears should perk up and you should anticipate being
able to apply at least some of the lessons in this module.

Finally, the term \textbf{endpoint} refers to a particular, addressable part
of an API that is normally defined by a base URL. For instance, the standard
Twitter search API is at \url{https://api.twitter.com/1.1/search/tweets.json}
and the endpoint for getting full information about specific tweets is \url{
https://api.twitter.com/1.1/statuses/lookup.json}. In general, an API will
have many different endpoints, each of which defines its own parameters and is
used for a particular type of information request.

\section{Accessing RESTful APIs with the \texttt{requests} package}

Python's \texttt{requests} package is an easy way to make a call to an
external API and get the results. It basically encapsulates the under-the-hood
stuff that formulates, sends, and parses the results of an HTTP request.

HTTP defines various operations such as \texttt{GET}, \texttt{POST}, and
\texttt{HEAD}. The \texttt{requests} package bundles these up in function
names with similar names: \texttt{requests.get()}, \texttt{requests.post()},
\textit{etc.} The most common method we'll use in API access is the
\texttt{.get()} command. It can actually be used to retrieve \textit{any}
Web-accessible page on the network:

\begin{Verbatim}[fontsize=\small,samepage=true,frame=none]
import requests
class_home_page = requests.get("http://cs.umw.edu/~stephen/data219")
print(class_home_page.content)
\end{Verbatim}

If you run the above code, you'll see the (gnarly) HTML code for the entire
course web page on your screen as one big string. Seeing this emphasizes that
\textit{when we execute RESTful API requests, we're really not doing much more
than accessing ``web pages," where ``web page" here means a set of results
that the server puts together for us.} This similarity to how browsers work is
what makes RESTful APIs so simple and powerful.

Most of the time we won't get HTML back, but rather a data format such as XML
or JSON. Predictably, I'll start with a Star Wars example. The SWAPI (Star
Wars API) is a neat little database that returns JSON data. Let's ask it for
``person with ID 13":

\small
\begin{Verbatim}[fontsize=\small,samepage=true,frame=none]
>>> response = requests.get("https://swapi.co/api/people/13")
>>> print(json.load(response.content))
{'name': 'Chewbacca', 'height': '228', 'mass': '112', 'hair_color': 'brown',
'skin_color': 'unknown', 'eye_color': 'blue', 'birth_year': '200BBY',
'gender': 'male', 'homeworld': 'https://swapi.co/api/planets/14/', 'films':
['https://swapi.co/api/films/2/', 'https://swapi.co/api/films/6/',
'https://swapi.co/api/films/3/', 'https://swapi.co/api/films/1/',
'https://swapi.co/api/films/7/'], 'species':
['https://swapi.co/api/species/3/'], 'vehicles':
['https://swapi.co/api/vehicles/19/'], 'starships':
['https://swapi.co/api/starships/10/', 'https://swapi.co/api/starships/22/'],
'created': '2014-12-10T16:42:45.066000Z', 'edited':
'2014-12-20T21:17:50.332000Z', 'url': 'https://swapi.co/api/people/13/'}
\end{Verbatim}
\normalsize

By calling \texttt{json.load()} on the \texttt{.content} of this response, we
get a Python dictionary with all Chewie's info. Note that some of the data
consists of \textit{other REST URLs}, a sign of a well-written API.

\small
\begin{Verbatim}[fontsize=\small,samepage=true,frame=none]
>>> chewie = json.load(response.content)
>>> print(json.loads(requests.get(chewie['homeworld']).content))
{'climate': 'tropical',
 'created': '2014-12-10T13:32:00.124000Z',
 'diameter': '12765',
 'edited': '2014-12-20T20:58:18.442000Z',
 'films': ['https://swapi.co/api/films/6/'],
 'gravity': '1 standard',
 'name': 'Kashyyyk',
 'orbital_period': '381',
 'population': '45000000',
 'residents': ['https://swapi.co/api/people/13/',
  'https://swapi.co/api/people/80/'],
 'rotation_period': '26',
 'surface_water': '60',
 'terrain': 'jungle, forests, lakes, rivers',
 'url': 'https://swapi.co/api/planets/14/'}
\end{Verbatim}
\normalsize



\section{The Twitter API}

Twitter's API uses the same principle, although it's of course larger and more
complicated. In addition, there's the hurdle of \textit{authenticating} your
Python program with the API so you can access the data.

\subsection{Authentication}

Many APIs require some form of \textbf{authentication}, which essentially
means you need to provide some personal information to identify you to the
site. Since we'll be dealing only with open APIs, this isn't a security thing
as much as it's a ``control of information" thing. By forcing you to register
with the site and provide some contact info, the organization retains a
measure of control and knowledge over how its information is used. Also, many
sites impose \textbf{rate limits} which constrain how much information you can
suck from the site every hour, and the way they track this is by associating
data requests with user accounts.

Every site has a different procedure for this, and the documentation varies
considerably in helpfulness. It's usually a matter of finding and connecting
the dots. Twitter, like many sites, uses the \textbf{OAuth} authentication
protocol (see \url{oauth.net}), an open standard that allows users to grant
one website permission to access data from another. (You've probably used this
without realizing it if you've ever enabled some third-party app which then
caused a confirmation dialog box to pop up saying, ``application so-and-so is
requesting access to the following information of yours...do you approve?")

For Twitter, here's the procedure I used:

\begin{enumerate}
\itemsep.1em
\item If you're not actually a Twitter user, create a Twitter account.
\item Sign on to \url{dev.twitter.com}.
\item Go to \url{apps.twitter.com} and create an ``application." For our
purposes, an application is sort of like a project, of which you can have
several. It's mainly an organizational thing.
\item Navigate to your application and choose ``Keys and Access Tokens."
Create an Access Token.
\end{enumerate}

You will now see, on the ``Keys and Access Tokens" page, \textit{four} strings
of gobbledygook, all four of which you'll need:

\vspace{.1in}
\small
\begin{compactitem}
\item A ``consumer key." Mine is \texttt{9dFucBtR1P58pJICXGV2gaKhj}.
\item A ``consumer secret." Mine is
\texttt{uKufTBegOrw5hrjNQmXFwcVtWjSb2bUHsMEoBwjN5gHzTbpJ0U}.
\item An ``access token." Mine is
\texttt{1019144197-5dhD0p1VbE80kPnzZrhjq43G08iMbguggkn67q1}.
\item An ``access token secret." Mine is
\texttt{JOtRHxgzhj6KQhuOuwljB214UImybVANFX6pbSYsDJhTd}.
\end{compactitem}
\normalsize
\vspace{.1in}

What do all these different things mean? Honestly, you don't need to know.
Just realize they're all necessary components of the way Twitter
authenticates, and if you just use them all in the right places, you'll be
able to get to any Twitter data you want, programmatically.

(Btw, the second and fourth of those items -- the ``secrets" -- are not
supposed to be shared with anyone, nor to appear in human-readable form in
your source code. For one-off data analysis projects, I confess I don't think
this is super important.)

\subsection{Authenticating with the \texttt{requests\_oauthlib} package}

To authenticate with OAuth, we'll use Python's \texttt{requests\_oauthlib}
package as follows:

\begin{Verbatim}[fontsize=\small,samepage=true,frame=none]
from requests_oauthlib import OAuth1
cons_key = "MlUmay5kA1vGWKokmmFofgRLX"
cons_secret = "2FFYCyI2rhUIEltkepeNhVcZvYufXJukCJMqE1s3ALKoLYm7LD"
access_token = "1019144197-cMFsHfxTZiG0oyOidzM7bCW4uX6PzjXVOyNQTUK"
access_secret = "fmW3KSw0kSYH43QPqSxUscnZD8XF8M8oEJnKm2sRAcq70"
oauth = OAuth1(cons_key, cons_secret, access_token, access_secret)
\end{Verbatim}

We now have an object called \texttt{oauth} which can be passed to functions
in the \texttt{requests} package to get data.

\subsection{HTTP GET parameters}

The Twitter API is documented at \url{https://dev.twitter.com/docs}, a page
well worth perusing in detail. You'll see numerous endpoints documented there,
each of which has its parameters listed when you click through.

Conceptually, making an API request is like making a function call: you're
triggering some operation, passing parameters, and getting a return value. The
syntax is a bit different though. In Python, you're used to making a function
call like this:

\begin{Verbatim}[fontsize=\small,samepage=true,frame=none]
get_user_timeline('katyperry',4)
\end{Verbatim}

to indicate something like ``get the latest four tweets from user
\texttt{@katyperry}." The corresponding syntax in URL Land is like this:

\footnotesize
\begin{Verbatim}[fontsize=\small,samepage=true,frame=none]
https://api.twitter.com/1.1/statuses/user_timeline.json?screen_name=katyperry&count=4
\end{Verbatim}
\normalsize

Note that parameters are ``passed" as key/value pairs, with no spaces, after a
``\texttt{?}", and separated by ``\texttt{\&}" signs. The following code:

\footnotesize
\begin{Verbatim}[fontsize=\small,samepage=true,frame=none]
four_from_katy = json.loads(requests.get(
    "https://api.twitter.com/1.1/statuses/user_timeline.json?screen_name=katyperry&count=4",
            auth=oauth).content)
for tweet in four_from_katy:
    print(tweet.text)
\end{Verbatim}
\normalsize

produces this result (today):

\begin{Verbatim}[fontsize=\small,samepage=true,frame=none]
You’re not going to miss this dance @effiepassero. #americanidol
@CadeFoehner Keep that feather, keep catching those dreams. #americanidol
HONEY!  #americanidol
@MaraJustine1 There’s not just glitter on your nails It’s in every pore,
and it comes out every time you sing  #americanidol
\end{Verbatim}

By the way, if a value you want to pass contains a space, you'll need to
convert that space to a plus sign (``\texttt{+}"). For instance, the standard
search endpoint is \texttt{https://api.twitter.com/1.1/search/tweets.json},
and requires a ``\texttt{q}" parameter with your query string. I could search
for Star Wars toys with:

\small
\begin{Verbatim}[fontsize=\small,samepage=true,frame=none]
toys = json.loads(requests.get(
    "https://api.twitter.com/1.1/search/tweets.json?q=star+wars+toys&count=5",
            auth=oauth).content)
\end{Verbatim}
\normalsize

This endpoint returns a dictionary of information:

\begin{Verbatim}[fontsize=\small,samepage=true,frame=none]
>>> print(list(toys.keys()))
['statuses', 'search_metadata']
\end{Verbatim}

Looks like the tweets themselves (``\texttt{statuses}") are in the
\texttt{statuses} key, so:

\small
\begin{Verbatim}[fontsize=\small,samepage=true,frame=none]
>>> for tweet in toys['statuses']:
        print(tweet['text'])
The #StarWars portfolio of #toys is about to get bigger. #Hasbro announced
it’s going to make #JabbatheHutt’s sail… https://t.co/FYC1DLftXX
RT @hottoysofficial: #HotToys 1/6th scale #ObiWanKenobi (Regular &amp; Deluxe
Version) collectible figures from #StarWars EP3 are available for…
HotToys 1/6 Old Luke Skywalker Star Wars Movie Series 12'' Man Figure
Collection https://t.co/NbHgtpnDwi #Toys https://t.co/ejrGXjVT3F
Star Wars Black Series  6" Rey (JEDI TRAINING )Toys R Us Exclusive
https://t.co/ArxT9kPjZt https://t.co/lcLbp7pwjG
\end{Verbatim}
\normalsize


If you have weird characters other than spaces, I suggest the \texttt{parse()}
method from the \texttt{urllib} library to make things kosher:

\begin{Verbatim}[fontsize=\small,samepage=true,frame=none]
import urllib.parse
>>> print(urllib.parse.quote_plus("http://cs.umw.edu/star wars toys! 20% off"))
'http%3A%2F%2Fcs.umw.edu%2Fstar+wars+toys%21+20%25+off'
\end{Verbatim}


\subsection{Specialized packages}

For many popular APIs, including Twitter, you may find Python packages written
to streamline and hide some of the protocol idiosyncrasies. These may be
easier to use, depending on how well written and maintained they are. For
Twitter, there are numerous ones: \texttt{python-twitter},
\texttt{TwitterAPI}, \texttt{tweepy}, \texttt{TweetPony}, \texttt{twython},
and more. This is a testimony to how powerful (and common) it is to be able to
access data through external APIs.

ADD SECTION FOR FOOD2FORK

\end{document}
